\documentclass[aspectratio=169]{beamer}
% talk-common.tex —— 各场分享共用的导言区配置
% 放在仓库根目录，由各个入口 .tex 通过 % talk-common.tex —— 各场分享共用的导言区配置
% 放在仓库根目录，由各个入口 .tex 通过 % talk-common.tex —— 各场分享共用的导言区配置
% 放在仓库根目录，由各个入口 .tex 通过 \input{talk-common.tex} 引用。
% 编译时需要把仓库根目录加入 TEXINPUTS，详见根目录 AGENTS.md 第 2 节。

\usetheme[customfont,showmaxslides,nofooterlogo]{pureminimalistic}
\usefonttheme{professionalfonts}
\usepackage[UTF8,fontset=none]{ctex}
\usepackage{graphicx}
\usepackage{listings}
\usepackage{appendixnumberbeamer}

\setsansfont{Fira Sans}
\setmonofont{Fira Mono}
\setCJKmainfont{Noto Sans CJK SC}
\setCJKsansfont{Noto Sans CJK SC}
\setCJKmonofont{Noto Sans CJK SC}

% 默认不引用上游示例 logo，有真实素材时再统一配置
\renewcommand{\logotitle}{}
\renewcommand{\logoheader}{}
\renewcommand{\logofooter}{}
\renewcommand{\pageword}{}

% 字号设置必须放在加载主题之后，避免被主题覆盖
\setbeamerfont{title}{size=\LARGE,shape=\upshape,series=\bfseries}
\setbeamerfont{presentation title}{shape=\upshape}
\setbeamerfont{frametitle}{size=\Large,shape=\upshape,series=\bfseries}

% 对浅色/深色都保持明确的代码对比度
\definecolor{TalkCodeBg}{HTML}{F3F4F6}
\definecolor{TalkCodeFg}{HTML}{17212B}
\definecolor{TalkCodeKeyword}{HTML}{174EA6}
\definecolor{TalkCodeComment}{HTML}{47634D}
\definecolor{TalkCodeString}{HTML}{8A3014}
\lstdefinestyle{talkcode}{
  basicstyle=\ttfamily\small\color{TalkCodeFg},
  backgroundcolor=\color{TalkCodeBg},
  keywordstyle=\bfseries\color{TalkCodeKeyword},
  commentstyle=\color{TalkCodeComment},
  stringstyle=\color{TalkCodeString},
  columns=fullflexible,
  keepspaces=true,
  showstringspaces=false,
  breaklines=true,
  tabsize=4,
  numbers=none,
  frame=none
}
\lstset{style=talkcode}
 引用。
% 编译时需要把仓库根目录加入 TEXINPUTS，详见根目录 AGENTS.md 第 2 节。

\usetheme[customfont,showmaxslides,nofooterlogo]{pureminimalistic}
\usefonttheme{professionalfonts}
\usepackage[UTF8,fontset=none]{ctex}
\usepackage{graphicx}
\usepackage{listings}
\usepackage{appendixnumberbeamer}

\setsansfont{Fira Sans}
\setmonofont{Fira Mono}
\setCJKmainfont{Noto Sans CJK SC}
\setCJKsansfont{Noto Sans CJK SC}
\setCJKmonofont{Noto Sans CJK SC}

% 默认不引用上游示例 logo，有真实素材时再统一配置
\renewcommand{\logotitle}{}
\renewcommand{\logoheader}{}
\renewcommand{\logofooter}{}
\renewcommand{\pageword}{}

% 字号设置必须放在加载主题之后，避免被主题覆盖
\setbeamerfont{title}{size=\LARGE,shape=\upshape,series=\bfseries}
\setbeamerfont{presentation title}{shape=\upshape}
\setbeamerfont{frametitle}{size=\Large,shape=\upshape,series=\bfseries}

% 对浅色/深色都保持明确的代码对比度
\definecolor{TalkCodeBg}{HTML}{F3F4F6}
\definecolor{TalkCodeFg}{HTML}{17212B}
\definecolor{TalkCodeKeyword}{HTML}{174EA6}
\definecolor{TalkCodeComment}{HTML}{47634D}
\definecolor{TalkCodeString}{HTML}{8A3014}
\lstdefinestyle{talkcode}{
  basicstyle=\ttfamily\small\color{TalkCodeFg},
  backgroundcolor=\color{TalkCodeBg},
  keywordstyle=\bfseries\color{TalkCodeKeyword},
  commentstyle=\color{TalkCodeComment},
  stringstyle=\color{TalkCodeString},
  columns=fullflexible,
  keepspaces=true,
  showstringspaces=false,
  breaklines=true,
  tabsize=4,
  numbers=none,
  frame=none
}
\lstset{style=talkcode}
 引用。
% 编译时需要把仓库根目录加入 TEXINPUTS，详见根目录 AGENTS.md 第 2 节。

\usetheme[customfont,showmaxslides,nofooterlogo]{pureminimalistic}
\usefonttheme{professionalfonts}
\usepackage[UTF8,fontset=none]{ctex}
\usepackage{graphicx}
\usepackage{listings}
\usepackage{appendixnumberbeamer}

\setsansfont{Fira Sans}
\setmonofont{Fira Mono}
\setCJKmainfont{Noto Sans CJK SC}
\setCJKsansfont{Noto Sans CJK SC}
\setCJKmonofont{Noto Sans CJK SC}

% 默认不引用上游示例 logo，有真实素材时再统一配置
\renewcommand{\logotitle}{}
\renewcommand{\logoheader}{}
\renewcommand{\logofooter}{}
\renewcommand{\pageword}{}

% 字号设置必须放在加载主题之后，避免被主题覆盖
\setbeamerfont{title}{size=\LARGE,shape=\upshape,series=\bfseries}
\setbeamerfont{presentation title}{shape=\upshape}
\setbeamerfont{frametitle}{size=\Large,shape=\upshape,series=\bfseries}

% 对浅色/深色都保持明确的代码对比度
\definecolor{TalkCodeBg}{HTML}{F3F4F6}
\definecolor{TalkCodeFg}{HTML}{17212B}
\definecolor{TalkCodeKeyword}{HTML}{174EA6}
\definecolor{TalkCodeComment}{HTML}{47634D}
\definecolor{TalkCodeString}{HTML}{8A3014}
\lstdefinestyle{talkcode}{
  basicstyle=\ttfamily\small\color{TalkCodeFg},
  backgroundcolor=\color{TalkCodeBg},
  keywordstyle=\bfseries\color{TalkCodeKeyword},
  commentstyle=\color{TalkCodeComment},
  stringstyle=\color{TalkCodeString},
  columns=fullflexible,
  keepspaces=true,
  showstringspaces=false,
  breaklines=true,
  tabsize=4,
  numbers=none,
  frame=none
}
\lstset{style=talkcode}


\title[技术部第二次分享]{技术部第二次技术分享会}
\subtitle{Java 基础语法 · 简明版}
\author{姚雨森}
\institute{技术部 分管主席}
\date{2026-10-23}

\begin{document}

\maketitle

% \begin{frame}[fragile]{输出：让程序主动“开口”}
%   怎样知道程序做了什么？先让它把信息打印到控制台。
%   \begin{lstlisting}[language=Java]
% public class Main {
%     public static void main(String[] args) {
%         // Print two lines.
%         System.out.println("Hello, Java!");
%         System.out.println("Let's go shopping!");
%     }
% }
%   \end{lstlisting}
%   \begin{itemize}
%     \item 示例从 \texttt{main} 开始，语句从上到下执行；每次输出后换行。
%     \item \texttt{//} 后的注释给人看，不作为语句执行。
%   \end{itemize}
%   \medskip
%   后面的代码片段都放在 \texttt{main} 的大括号内。
% \end{frame}

\begin{frame}[fragile]{让程序主动“开口”}
  我们需要知道程序正在执行什么，可以让它把信息打印到控制台。
  \par\medskip
  使用 \texttt{System.out.println(...)} 输出一行文字：
  \begin{lstlisting}[language=Java]
System.out.println("Hello, Java!");
  \end{lstlisting}
  双引号里面的文字会被打印出来，双引号本身不会显示。
  \par\medskip
  控制台输出：
  \begin{lstlisting}[language={}]
Hello, Java!
  \end{lstlisting}
\end{frame}

\begin{frame}[fragile]{把输出语句写在 main 中}
  程序从 \texttt{main} 方法开始执行。先把输出语句写在它的大括号内：
  \begin{lstlisting}[language=Java]
public class Main {
    public static void main(String[] args) {
        System.out.println("First");
        System.out.println("Second");
    }
}
  \end{lstlisting}
  这里的两条语句\textbf{从上到下依次执行}，每次输出后换行：
  \begin{lstlisting}[language={}]
First
Second
  \end{lstlisting}
\end{frame}

\begin{frame}[fragile]{注释：给人看的说明}
  注释用于解释代码，注释中的内容不会被当作语句执行。
  \begin{itemize}
    \item 单行注释：\texttt{//} 后面直到本行末尾的内容。
    \item 多行注释：\texttt{/*} 与 \texttt{*/} 之间的内容，可以跨行。
  \end{itemize}
  \medskip
  下面的代码写在 \texttt{main} 中：
  \begin{lstlisting}[language=Java]
// Print a greeting.
System.out.println("Hello!"); // This is a comment.
/* The next statement is commented out.
System.out.println("Hidden");
*/
System.out.println("Goodbye!");
  \end{lstlisting}
  控制台依次输出 \texttt{Hello!} 和 \texttt{Goodbye!}，各占一行；
  \texttt{Hidden} 不会输出。
\end{frame}

% \begin{frame}[fragile]{变量：把需要的数据记下来}
%   商品数量可能变化，我们需要保存它，也需要修改它。
%   \begin{lstlisting}[language=Java]
% int quantity = 2;
% System.out.println(quantity); // 2

% quantity = quantity + 1;
% System.out.println(quantity); // 3
%   \end{lstlisting}
%   \begin{itemize}
%     \item \texttt{int quantity = 2;}：声明整数变量，并给它一个初始值。
%     \item \texttt{=}：先计算右边，再把结果保存到左边的变量中。
%     \item 读取或重新赋值时，直接使用变量名，不再写类型。
%   \end{itemize}
% \end{frame}

\begin{frame}[fragile]{变量：保存程序需要的数据}
  程序有时需要做运算，也需要保存一些数据，比如数值、文本。
  \par\medskip
  可以用\textbf{变量}保存这些数据，再通过变量名读取或修改它们。
%   \par\medskip
%   下面的代码写在 \texttt{main} 中：
%   \begin{lstlisting}[language=Java]
% int price = 60;
% System.out.println(price);

% price = 70;
% System.out.println(price);
%   \end{lstlisting}
%   \begin{itemize}
%     \item \texttt{int price = 60;}：声明整数变量 \texttt{price}，初始化为 60。
%     \item \texttt{price = 70;}：重新赋值，把 \texttt{price} 中保存的值改为 70。
%     \item 两次 \texttt{println} 依次输出 60 和 70，各占一行。
%   \end{itemize}
\end{frame}

\begin{frame}[fragile]{变量的声明与赋值}
  \textbf{声明变量时赋值（初始化）}
  \par
  语法：\quad 类型\quad 变量名\quad \texttt{=}\quad 初始值\texttt{;}
  \begin{lstlisting}[language=Java]
int price = 60;
  \end{lstlisting}
  \medskip
  \textbf{也可以先声明，再赋值}
  \par
  声明：\quad 类型\quad 变量名\texttt{;}
  \par
  赋值：\quad 变量名\quad \texttt{=}\quad 值\texttt{;}
  \begin{lstlisting}[language=Java]
int price;
price = 60;
System.out.println(price);
  \end{lstlisting}
  赋值时不用再写类型。
\end{frame}

\begin{frame}{数据类型：不同信息有不同表达}
  商品数量、单价、是否有货和商品名，需要表达不同性质的信息。
  \medskip
  \begin{center}
    \renewcommand{\arraystretch}{1.45}
    \begin{tabular}{lll}
      \textbf{要保存什么} & \textbf{常用类型} & \textbf{值的例子} \\
      \hline
      整数数量 & \texttt{int} & \texttt{2} \\
      带小数的数值 & \texttt{double} & \texttt{12.5} \\
      是或否 & \texttt{boolean} & \texttt{true} \\
      单个字符 & \texttt{char} & \texttt{'A'} \\
      一段文本 & \texttt{String} & \texttt{"Notebook"} \\
    \end{tabular}
  \end{center}
  \medskip
  \alert{先想清楚数据是什么，再选择类型。}
\end{frame}

% \begin{frame}[fragile]{数字：整数与小数怎么选}
%   数量通常是整数，测量值、单价等数值可能带小数。
%   \begin{lstlisting}[language=Java]
% int quantity = 2;
% double price = 12.5;
% long totalVisits = 3000000000L;
%   \end{lstlisting}
%   \begin{itemize}
%     \item 常见整数用 \texttt{int}；需要更大范围时可用 \texttt{long}。
%     \item 小数通常用 \texttt{double}；浮点数可能是近似值。
%     \item 还有 \texttt{byte}、\texttt{short}、\texttt{float}，按具体需求再了解。
%   \end{itemize}
%   \medskip
%   类型能表示的范围和精度有限，不需要现在背下全部数值。
% \end{frame}

\begin{frame}[fragile]{boolean：表达“是或否”}
  商品有没有货？用户是不是会员？这些状态只有两种答案。
  \begin{lstlisting}[language=Java]
boolean inStock = true;
boolean isMember = false;

System.out.println(inStock);  // true
System.out.println(isMember); // false
  \end{lstlisting}
  \begin{itemize}
    \item \texttt{true} 表示真，\texttt{false} 表示假。
    \item \texttt{boolean} 只有这两个值，不用数字 0 和 1 代替。
  \end{itemize}
  \medskip
  后面可以根据这些状态，让程序采取不同的动作。
\end{frame}

\begin{frame}[fragile]{字符与字符串：保存文字}
  一个等级字母可以用 \texttt{char}，商品名称可以用 \texttt{String}。
  \begin{lstlisting}[language=Java]
char grade = 'A';
String name = "Notebook";

System.out.println(name);   // Notebook
System.out.println("name"); // name
  \end{lstlisting}
  \begin{itemize}
    \item 字符用单引号，如 \texttt{'A'}；文本用双引号，如 \texttt{"Notebook"}。
    \item \texttt{name} 读取变量的值，\texttt{"name"} 表示文字本身。
    \item \texttt{String} 首字母大写，不属于基本数据类型。
  \end{itemize}
\end{frame}

% \begin{frame}[fragile]{运算：从单价和数量得到总价}
%   买两种商品：分别计算金额，再把结果加起来。
%   \begin{lstlisting}[language=Java]
% double priceA = 12.5;
% double priceB = 8.0;
% int quantityA = 2;
% int quantityB = 3;
% double total = priceA * quantityA + priceB * quantityB;
% System.out.println(total); // 49.0
%   \end{lstlisting}
%   \begin{itemize}
%     \item \texttt{+ - * /} 表示加减乘除，\texttt{\%} 求余数；括号中的部分先算。
%     \item 整数相除：\texttt{7 / 2} 得到 3；\texttt{7.0 / 2} 得到 3.5。
%     \item 也可先转换类型：\texttt{(double) 7 / 2} 得到 3.5。
%   \end{itemize}
% \end{frame}

\begin{frame}[fragile]{字符串拼接：让结果看得懂}
  只看到一个数字还不够，我们可以给结果加上说明。
  \begin{lstlisting}[language=Java]
String name = "Notebook";
double total = 25.0;
System.out.println(name + " total: " + total);
  \end{lstlisting}
  输出：\quad \texttt{Notebook total: 25.0}
%   \par\bigskip
%   \texttt{+} 遇到字符串时进行拼接。想先计算，就加括号：
%   \begin{lstlisting}[language=Java]
% System.out.println("Count: " + (2 + 3)); // Count: 5
  % \end{lstlisting}
\end{frame}

\begin{frame}[fragile]{算术运算符：加、减、乘、除与取余}
  我们可以对数字做运算，把结果保存到变量中，或直接打印出来。
  \begin{center}
    \renewcommand{\arraystretch}{1.15}
    \begin{tabular}{cllc}
      \textbf{符号} & \textbf{运算} & \textbf{表达式示例} & \textbf{结果} \\
      \hline
      \texttt{+} & 加法 & \texttt{7 + 2} & 9 \\
      \texttt{-} & 减法 & \texttt{7 - 2} & 5 \\
      \texttt{*} & 乘法 & \texttt{7 * 2} & 14 \\
      \texttt{/} & 除法 & \texttt{7 / 2} & 3 \\
      \texttt{\%} & 取余 & \texttt{7 \% 2} & 1 \\
    \end{tabular}
  \end{center}
  \begin{lstlisting}[language=Java]
int sum = 7 + 2;
System.out.println(sum);   // 9
  \end{lstlisting}
  \texttt{\%} 求除法的余数，常被称为“取模”；这里先看正整数。
\end{frame}

% \begin{frame}[fragile]{条件：把规则变成真假判断}
%   “总价超过 40 元，并且是会员”，可以写成条件表达式。
%   \begin{lstlisting}[language=Java]
% double total = 49.0;
% boolean isMember = true;
% boolean eligible = total > 40 && isMember;
% System.out.println(eligible); // true
%   \end{lstlisting}
%   \begin{itemize}
%     \item 比较：\texttt{> >= < <=}，以及相等 \texttt{==}、不等 \texttt{!=}。
%     \item 组合：\texttt{\&\&} 表示“并且”，\texttt{||} 表示“或者”，\texttt{!} 表示“取反”。
%     \item 判断的结果是 \texttt{true} 或 \texttt{false}。
%   \end{itemize}
%   \medskip
%   注意：\texttt{=} 是赋值；\texttt{==} 才是相等判断。
% \end{frame}

\begin{frame}[fragile]{比较运算：结果是 boolean}
  比较两个数值，结果是 \texttt{true} 或 \texttt{false}，可以保存到 \texttt{boolean} 变量中。
  \begin{center}
    \renewcommand{\arraystretch}{1.05}
    \begin{tabular}{cllc}
      \textbf{运算符} & \textbf{含义} & \textbf{示例} & \textbf{结果} \\
      \hline
      \texttt{==} & 等于 & \texttt{7 == 2} & \texttt{false} \\
      \texttt{!=} & 不等于 & \texttt{7 != 2} & \texttt{true} \\
      \texttt{>} & 大于 & \texttt{7 > 2} & \texttt{true} \\
      \texttt{<} & 小于 & \texttt{7 < 2} & \texttt{false} \\
      \texttt{>=} & 大于或等于 & \texttt{7 >= 7} & \texttt{true} \\
      \texttt{<=} & 小于或等于 & \texttt{7 <= 2} & \texttt{false} \\
    \end{tabular}
  \end{center}
  \begin{lstlisting}[language=Java]
boolean passed = 80 >= 60;
System.out.println(passed); // true
  \end{lstlisting}
  注意：\texttt{=} 是赋值，\texttt{==} 才是判断是否相等。
\end{frame}

\begin{frame}[fragile]{逻辑运算：组合或取反条件}
  逻辑运算的操作数和结果都是 \texttt{boolean}。
  \begin{center}
    \renewcommand{\arraystretch}{1.25}
    \begin{tabular}{cll}
      \textbf{运算符} & \textbf{含义} & \textbf{结果为 true 的条件} \\
      \hline
      \texttt{\&\&} & 与 & 两边都为 \texttt{true} \\
      \texttt{||} & 或 & 至少一边为 \texttt{true} \\
      \texttt{!} & 非（取反） & 原来的值为 \texttt{false} \\
    \end{tabular}
  \end{center}
  \begin{lstlisting}[language=Java]
int score = 80;
boolean inRange = (score >= 0) && (score <= 100); // true
boolean outside = (score < 0) || (score > 100);   // false
boolean failed = !(score >= 60);                // false
  \end{lstlisting}
\end{frame}

\begin{frame}{综合示例：商品总价计算需求}
  购买 A、B 两种商品，根据各自的单价和数量计算总价。
  \par\bigskip
  \begin{itemize}
    \item 用 \texttt{priceA}、\texttt{priceB} 保存两种商品的单价，单位为元。
    \item 用 \texttt{quantityA}、\texttt{quantityB} 保存各自的购买数量。
    \item 直接修改代码中四个变量的初始值。
    \item 计算总价，并在控制台打印结果。
  \end{itemize}
  \bigskip
  \textbf{计算公式：}总价 = A 的单价 $\times$ 数量 + B 的单价 $\times$ 数量。
  \par\medskip
  示例数据：A 每件 12.5 元，买 2 件；B 每件 8 元，买 3 件。
\end{frame}

\begin{frame}[fragile]{综合示例：商品总价计算实现}
  \begin{lstlisting}[language=Java]
public class Main {
    public static void main(String[] args) {
        double priceA = 12.5;
        double priceB = 8.0;
        int quantityA = 2;
        int quantityB = 3;
        double total = priceA * quantityA + priceB * quantityB;
        System.out.println("Total: " + total);
    }
}
  \end{lstlisting}
  \begin{itemize}
    \item 先分别计算两种商品的金额，再相加得到总价。
    \item 本例输出 \texttt{Total: 49.0}，即总价为 49 元。
  \end{itemize}
\end{frame}

\begin{frame}{新需求：满足条件时才打折}
  现在增加一个需求：总价\textbf{超过 40 元}时打八折。
  \par
  也就是：当 \texttt{total > 40} 时，让 \texttt{total} 乘以 0.8；否则保持不变。
\end{frame}

\begin{frame}[fragile]{if：满足条件时才执行}
  \textbf{if 语法：条件为 true 时执行，否则跳过。}
  \begin{lstlisting}[language=Java,escapeinside={(*@}{@*)}]
if ((*@条件@*)) {
    (*@条件成立时执行的语句@*)
}
  \end{lstlisting}
  条件必须是 \texttt{boolean}。例如：总价超过 40 元时打八折。
  \begin{lstlisting}[language=Java]
double total = 49.0;
if (total > 40) {
    total = total * 0.8;
}
System.out.println("Total: " + total); // Total: 39.2
  \end{lstlisting}
\end{frame}

\begin{frame}{新需求：打折与未打折时输出不同文字}
  如果打折了，输出文字也要变化：
  \begin{itemize}
    \item 打折时，输出 \texttt{Discounted total:} 和折后总价。
    \item 未打折时，输出 \texttt{Total:} 和正常总价。
  \end{itemize}
  \bigskip
  原来的写法只在 \texttt{if} 中修改价格，最后总会执行外面的同一条输出语句，
  因而无法区分这两种文字。
  \par\bigskip
  自然的想法是：\textbf{如果打折了，就按折扣文本打印；否则，按正常文本打印。}
  \par\medskip
  用 \texttt{else} 表示“否则”，让程序在两条分支中选择一条执行。
\end{frame}

\begin{frame}[fragile]{if--else：在两条分支中选择一条}
  条件为 \texttt{true} 时执行 \texttt{if} 分支，为 \texttt{false} 时执行 \texttt{else} 分支。
  \par\medskip
  \begin{columns}[T,onlytextwidth]
    \begin{column}{0.35\textwidth}
      \textbf{语法}
  \begin{lstlisting}[language=Java,escapeinside={(*@}{@*)}]
if ((*@条件@*)) {
    (*@条件成立时执行@*)
} else {
    (*@条件不成立时执行@*)
}
  \end{lstlisting}
    \end{column}
    \begin{column}{0.62\textwidth}
      \textbf{折扣示例}
  \begin{lstlisting}[language=Java]
double total = 49.0;
if (total > 40) {
    total = total * 0.8;
    System.out.println(
        "Discounted total: " + total);
} else {
    System.out.println(
        "Total: " + total);
}
  \end{lstlisting}
    \end{column}
  \end{columns}
\end{frame}

\begin{frame}{新需求：按总价划分价格档次}
  假设总价非负，我们希望根据金额输出对应的价格档次。
  \par\bigskip
  \begin{center}
    \renewcommand{\arraystretch}{1.4}
    \begin{tabular}{ll}
      \textbf{总价范围} & \textbf{档次} \\
      \hline
      $0 \leq \texttt{total} \leq 40$ & 便宜（\texttt{Cheap}） \\
      $40 < \texttt{total} \leq 100$ & 适中（\texttt{Moderate}） \\
      $\texttt{total} > 100$ & 昂贵（\texttt{Expensive}） \\
    \end{tabular}
  \end{center}
  \bigskip
  40 元属于“便宜”，100 元属于“适中”，每个总价只属于一个档次。
\end{frame}

\begin{frame}[fragile]{先用三个独立的 if 实现}
  \begin{lstlisting}[language=Java]
double total = 80.0;
if (total <= 40) {
    System.out.println("Cheap");
}
if (total > 40 && total <= 100) {
    System.out.println("Moderate");
}
if (total > 100) {
    System.out.println("Expensive");
}
  \end{lstlisting}
  这三个判断在逻辑上是一组：对\textbf{同一个不变的总价}，恰好只有一个条件成立。
  \par\medskip
  但三个独立的 \texttt{if} 仍会从上到下，逐个判断条件。
\end{frame}

\begin{frame}[fragile]{修改总价后，后面的 if 可能再次成立}
  假设超过 100 元打五折，并把这个判断放在另外两个判断之间：
  \begin{lstlisting}[language=Java]
double total = 180.0;
if (total <= 40) {
    System.out.println("Cheap");
}
if (total > 100) {
    total = total * 0.5;
    System.out.println("Discounted");
}
if (total > 40 && total <= 100) {
    System.out.println("Moderate");
}
  \end{lstlisting}
  180 元变成 90 元后，又满足“适中”条件，两个分支都执行了。
\end{frame}

\begin{frame}[fragile]{else if：一组条件，只执行一个分支}
  我们需要把多个判断连成一组：\textbf{一旦某个条件成立，后面的条件就不再判断。}
  \par\medskip
  \begin{lstlisting}[language=Java,escapeinside={(*@}{@*)}]
if ((*@条件 1@*)) {
    (*@条件 1 成立时执行的语句@*)
} else if ((*@条件 2@*)) {
    (*@条件 2 成立时执行的语句@*)
} else {
    (*@以上条件都不成立时执行的语句@*)
}
  \end{lstlisting}
  \begin{itemize}
    \item 从上到下判断，只执行\textbf{第一个成立条件}对应的分支。
    \item 即使分支中修改了变量，也不会继续判断后面的 \texttt{else if}。
    \item 最后的 \texttt{else} 可以省略；省略后可能一个分支也不执行。
  \end{itemize}
\end{frame}

\begin{frame}[fragile]{新需求：余额不足时充值}
  商品总价 180 元，余额只有 50 元。假设每次充值增加 50 元。
  \par\medskip
  先用 \texttt{if}：余额不足，就充值一次。
  \begin{lstlisting}[language=Java]
double total = 180.0;
double balance = 50.0;
if (balance < total) {
    balance = balance + 50;
}
System.out.println("Balance: " + balance); // 100.0
  \end{lstlisting}
  \textbf{问题：充值一次后，余额仍然不够付款。}
  \par\medskip
  \texttt{if} 只判断一次，不会自动回头检查余额是否已经足够。
\end{frame}

\begin{frame}{充值后重新检查，直到余额足够}
  我们需要的是：\textbf{只要余额不足，就继续充值。}
  \par\bigskip
  \begin{enumerate}
    \item 检查余额是否小于商品总价。
    \item 如果不足，充值 50 元，然后回到第一步。
    \item 如果已经足够，结束充值，继续付款。
  \end{enumerate}
  \bigskip
  本例余额变化：$50 \rightarrow 100 \rightarrow 150 \rightarrow 200$ 元。
  \par\medskip
  余额达到 200 元时，不再满足“余额不足”，于是停止充值。
  \par\medskip
  用 \texttt{while} 表示这种\textbf{条件成立就反复执行}的过程。
\end{frame}

\begin{frame}[fragile]{while：条件成立就重复执行}
  每轮先判断条件：为 \texttt{true} 时执行循环体，再回到条件处；为 \texttt{false} 时退出。
  \par\medskip
  \begin{columns}[T,onlytextwidth]
    \begin{column}{0.33\textwidth}
      \textbf{语法}
      \begin{lstlisting}[language=Java,escapeinside={(*@}{@*)}]
while ((*@条件@*)) {
    (*@反复执行的语句@*)
}
      \end{lstlisting}
      条件必须是 \texttt{boolean}。
    \end{column}
    \begin{column}{0.64\textwidth}
      \textbf{充值与付款}
      \begin{lstlisting}[language=Java]
double total = 180.0;
double balance = 50.0;
while (balance < total) {
    balance = balance + 50;
}
balance = balance - total;
System.out.println(
    "Remaining: " + balance); // 20.0
      \end{lstlisting}
    \end{column}
  \end{columns}
  \medskip
  \begin{itemize}
    \item 一开始余额就足够时，循环一次也不执行，直接付款。
    \item 如果余额不足，却不在循环中增加余额，就会一直循环。
  \end{itemize}
\end{frame}

\begin{frame}[fragile]{用 while 计数：打印商品编号}
  新需求：依次打印第 1～5 件商品的编号。
  \begin{lstlisting}[language=Java]
int number = 1;
while (number <= 5) {
    System.out.println("Item: " + number);
    number++;
}
  \end{lstlisting}
  \begin{itemize}
    \item 初始化：\texttt{int number = 1;}，从 1 开始。
    \item 继续条件：\texttt{number <= 5}，编号不超过 5。
    \item 更新：\texttt{number++}，每轮结束后加 1。
  \end{itemize}
  \medskip
  三个部分分散在不同位置；如果忘记 \texttt{number++}，就会一直打印编号 1。
\end{frame}

\begin{frame}[fragile]{for：把计数循环的三个部分写在一起}
  \texttt{for} 把初始化、继续条件和更新集中到循环开头，便于检查。
  \begin{lstlisting}[language=Java,escapeinside={(*@}{@*)}]
for ((*@初始化@*); (*@继续条件@*); (*@更新@*)) {
    (*@每轮执行的语句@*)
}
  \end{lstlisting}
  同样打印第 1～5 件商品的编号：
  \begin{lstlisting}[language=Java]
for (int number = 1; number <= 5; number++) {
    System.out.println("Item: " + number);
}
  \end{lstlisting}
  \begin{itemize}
    \item 初始化只执行一次；每轮先判断条件，再执行循环体。
    \item 循环体结束后执行更新，再回到条件处；条件为假时结束。
    \item \texttt{for} 仍需写更新语句，并不会自动帮我们加 1。
  \end{itemize}
\end{frame}

% \begin{frame}[fragile]{if / else：根据情况选择动作}
%   新需求：总价超过 40 元就打八折，否则保持原价。
%   \begin{lstlisting}[language=Java]
% double total = 49.0;
% if (total > 40) {
%     total = total * 0.8;
%     System.out.println("Discount applied");
% } else {
%     System.out.println("No discount");
% }
% System.out.println("Total: " + total);
%   \end{lstlisting}
%   \begin{itemize}
%     \item 条件成立执行 \texttt{if} 分支，否则执行 \texttt{else} 分支。
%     \item 大括号圈出同一分支的操作；只需“满足才做”时，可以省略 \texttt{else}。
%   \end{itemize}
% \end{frame}

% \begin{frame}[fragile]{else if：多个情况选择一种}
%   新需求：根据总价，把订单分成三个档次。
%   \begin{lstlisting}[language=Java]
% double total = 80.0;
% if (total > 100) {
%     System.out.println("Large order");
% } else if (total > 40) {
%     System.out.println("Medium order");
% } else {
%     System.out.println("Small order");
% }
%   \end{lstlisting}
%   \begin{itemize}
%     \item 从上到下判断，只执行\textbf{第一个成立条件}对应的分支。
%     \item 本例输出 \texttt{Medium order}；所有条件都不成立时执行 \texttt{else}。
%   \end{itemize}
% \end{frame}

% \begin{frame}[fragile]{while：一次不够，就继续做}
%   模拟余额不足时每次增加 50 元，直到足够支付 180 元。
%   \begin{lstlisting}[language=Java]
% int balance = 50;
% int total = 180;
% while (balance < total) {
%     balance = balance + 50;
% }
% System.out.println(balance); // 200
%   \end{lstlisting}
%   \begin{center}
%     余额：\quad $50 \rightarrow 100 \rightarrow 150 \rightarrow 200$ \quad 停止
%   \end{center}
%   \begin{itemize}
%     \item 每轮先判断；条件成立就执行，再回到条件处。
%     \item 一开始就足够时，一次也不执行；余额始终不变则可能一直循环。
%   \end{itemize}
% \end{frame}

% \begin{frame}[fragile]{for：有规律地重复}
%   新需求：依次打印第 1～5 件商品的编号。
%   \begin{lstlisting}[language=Java]
% for (int number = 1; number <= 5; number++) {
%     System.out.println("Item: " + number);
% }
%   \end{lstlisting}
%   \begin{enumerate}
%     \item 起点：\texttt{int number = 1}，只在开始时执行一次。
%     \item 条件：\texttt{number <= 5}，每轮开始前检查。
%     \item 更新：\texttt{number++}，每轮执行完后将编号加一。
%   \end{enumerate}
%   \medskip
%   \alert{先初始化，再反复执行“判断 → 循环体 → 更新”。}
% \end{frame}

% \begin{frame}{分支与循环：什么时候用哪一个}
%   \begin{center}
%     \renewcommand{\arraystretch}{1.65}
%     \begin{tabular}{lll}
%       \textbf{需求} & \textbf{常用写法} & \textbf{购物中的例子} \\
%       \hline
%       根据条件选择动作 & \texttt{if / else} & 满足条件才打折 \\
%       多个情况选择一种 & \texttt{else if} & 按金额划分档次 \\
%       条件成立就继续 & \texttt{while} & 余额不足就继续增加 \\
%       按计数规律重复 & \texttt{for} & 打印第 1～5 件商品编号 \\
%     \end{tabular}
%   \end{center}
%   \bigskip
%   \texttt{while} 和 \texttt{for} 常能互相改写，选择能把过程表达清楚的写法。
%   \par\medskip
%   \alert{分支决定“做哪一件”，循环决定“是否继续做”。}
% \end{frame}

% \begin{frame}[fragile]{循环 + 判断：筛选符合要求的数字}
%   小练习：打印 1～1000 中能被 7 整除的数。
%   \begin{lstlisting}[language=Java]
% for (int number = 1; number <= 1000; number++) {
%     if (number % 7 == 0) {
%         System.out.println(number);
%     }
% }
%   \end{lstlisting}
%   \begin{itemize}
%     \item \texttt{for}：依次检查 1～1000。
%     \item \texttt{number \% 7 == 0}：除以 7 的余数为 0。
%     \item \texttt{if}：只打印符合条件的数。
%   \end{itemize}
%   \medskip
%   输出：7、14、21、……、994，每个数占一行。
% \end{frame}

% \begin{frame}[fragile]{再加一个变量，就能统计数量}
%   新需求：除了打印，还要知道一共有多少个。
%   \begin{lstlisting}[language=Java]
% int count = 0;
% for (int number = 1; number <= 1000; number++) {
%     if (number % 7 == 0) {
%         System.out.println(number);
%         count++;
%     }
% }
% System.out.println("Count: " + count); // Count: 142
%   \end{lstlisting}
%   \begin{itemize}
%     \item \texttt{count} 在循环前初始化，每找到一个就加一。
%     \item 如果每轮都把它重设为 0，就无法保留前面累计的结果。
%   \end{itemize}
% \end{frame}

\begin{frame}{遇到问题时，想到这些工具}
  \begin{center}
    \renewcommand{\arraystretch}{1.4}
    \begin{tabular}{ll}
      \textbf{我想做什么} & \textbf{可以用什么} \\
      \hline
      看见程序的结果 & 输出 \\
      保存和更新数据 & 变量 \\
      表达不同性质的信息 & 数据类型 \\
      计算数值、组织文字 & 运算、字符串拼接 \\
      根据情况执行不同操作 & 分支 \\
      重复执行一段操作 & 循环 \\
    \end{tabular}
  \end{center}
  \bigskip
  \alert{先想清楚数据、条件和重复过程，再查具体写法。}
\end{frame}

\begin{frame}{接下来需要自主学习的内容}
  \begin{itemize}
    \item \textbf{（必选）}Java 基础语法完整版 PPT
      \bigskip
    \item \textbf{（可选）}《Java 核心技术 卷 I》第三章
  \end{itemize}
\end{frame}

\begin{frame}{课后练习}
  \begin{enumerate}
    \item \textbf{（必选）}完成补充语句的练习项目。
  \end{enumerate}
  \bigskip
  项目链接之后会发到群里。
\end{frame}

\end{document}
